\section{Discussion}
The reduction separates three questions that a single lifetime or valley count cannot answer. First, does a sector have a direct Hall weight? Second, which physical distributions does its collision coupling excite in the active sector? Third, how accurately must a particular readout be reproduced? A zero answer to the first does not settle the others.

The ablation assigns conditional roles rather than a universal ranking. Active valley distortion carries the retained redistribution coordinate, but only the coupled solution fixes its amplitude. Passive current is indispensable for ordinary conduction and can also alter Hall response. Passive valley polarization closes an additional finite-range feedback path despite having neither direct current nor Hall weight in the massless limit. It is redundant for the driven contact problem and can be unnecessary for a weak-coupling Hall-only tolerance. Four modes provide a common accurate description on the tested domain, not a theorem that every four-valley problem has four relevant modes.

The leakage calculation refines that statement. Contact closure follows from the Bloch-moment structure, while finite range produces a genuinely nonzero invariant residual. Hall accuracy is better than distribution accuracy because projection error is sampled by an observable-specific dual. Compensated drive and leakage forcing, small omitted output sensitivity, and signed modal residues explain the representative errors; uniformly fast leakage does not. Near a tiny net correction, even a small error relative to $|H_0|$ can be a large fraction of that correction. No precision for such a cancellation residual follows from a small total-response error.

These results adapt established sector-reduction and output-error concepts to a specified multivalley anomalous-velocity omission problem. Schur/Kron reduction and adjoint estimation are standard~\cite{DorflerBullo2013,PierceGiles2004}. Valley-polarizing relaxons and scattering-channel subtraction already demonstrate the kinetic content of nonlinear Hall response~\cite{LihmPark2024}. Our contribution is the physically defined ablation, analytic contact closure, quantified finite-range leakage and tolerance-specific distinction between scalar, redistribution and fixed-budget cancellation regimes. Neither the mathematics nor the general importance of valley relaxation is claimed as a priority result.

The physical scope is limited. Side-jump velocities, anomalous distribution terms and skew scattering are excluded and can occur at the same disorder order as the retained channel~\cite{NandySodemann2019,Du2019}. Gaussian disorder removes the third cumulant, not every such contribution; reducing $\Gamma_0$ does not establish a complete Hall response. Kernel and relative-texture controls constrain the present assumptions without demonstrating disorder independence. The bands and disorder ensemble are phenomenological, and the Fermi-contour calculation is at zero temperature. Finite temperature, boundaries and spatial ballistic-to-diffusive transport~\cite{Liu2026} require different calculations. No quantitative SnTe prediction, universal omission threshold or universal enhancement sign follows.
